% Part I: how to build, install and run every part of nabia.

\chapter{Getting started}\label{ch:guide}

\section{What is in the repository}
\begin{center}\small
\begin{tabularx}{\linewidth}{@{}l l X@{}}
\toprule
Folder & Open in & Contents \\
\midrule
\code{cpp/} & CLion (CMake) & the C++ library (\code{include}, \code{src}), Catch2 tests, the command-line
  backtest, and the pybind11 bindings that build the \code{nabia\_python\_bridge} wheel \\
\code{nabia\_python/} & PyCharm & the Python project: the \code{nabia} package (portfolio backtests, analytics,
  market data, and the sentiment risk-budgeting pipeline \code{nabia.sentiment\_rb}), its tests, its single
  \code{pyproject.toml} and its \code{.venv} \\
\code{doc/} & --- & this document (\code{nabia.tex}, chapters, the snippet extractor, \code{build.sh}) and the
  specification of the sentiment pipeline (\code{sentiment\_risk\_budgeting\_spec.md}) \\
\code{dist/} & --- & built wheels (not versioned) \\
\code{.github/workflows/wheels.yml} & --- & wheels for every platform in GitHub Actions \\
\bottomrule
\end{tabularx}
\end{center}
The two projects meet only through the wheel: \code{cpp/} builds \code{nabia\_python\_bridge} into
\code{dist/}, and \code{nabia\_python/} installs it into its \code{.venv}. The order of the steps below follows
from that.

\section{Prerequisites}
\begin{itemize}
  \item \textbf{C++}: CMake 3.25 or newer and a C++20 compiler (Apple clang from Xcode, GCC 12+, MSVC 2022).
        Eigen, iir1, Catch2, nlohmann/json and pybind11 are found on the system or downloaded by CMake.
  \item \textbf{Python}: CPython 3.10--3.14, and \href{https://docs.astral.sh/uv/}{uv} to build wheels for several
        versions. pdflatex (TeX Live, MacTeX) to build this document.
  \item \textbf{Optional}: the GitHub CLI (\code{gh}, logged in) to download published wheels or start CI; an
        Anthropic API key for the pipeline's \code{--live} mode.
\end{itemize}

\section{Step 1: the C++ library}
\begin{lstlisting}[language=bash]
cd cpp
cmake --preset debug                    # configure (once, or after editing a CMakeLists.txt)
cmake --build --preset debug -j         # library, tests, command-line application, C++ examples
ctest --preset debug                    # the Catch2 suite and an end-to-end run of the application
\end{lstlisting}
Presets: \code{debug}, \code{release}, \code{releaseWithDebugInfo}, \code{asan} (address and undefined-behaviour
sanitizers) and \code{python-bridge} (a development build of the bindings, \S\ref{sec:guide-debug}). In CLion,
open \code{cpp/}; it reads \code{CMakePresets.json}: enable the \code{debug} preset and run or debug
\code{nabia\_tests} or \code{nabia\_rb\_backtest}. The command-line backtest (Chapter~\ref{c:sec:cli}):
\begin{lstlisting}[language=bash]
./build/debug/apps/rb_backtest/nabia_rb_backtest --prices apps/rb_backtest/example/prices.csv \
    --config apps/rb_backtest/example/portfolio.json --output /tmp/result.csv
\end{lstlisting}

\section{Step 2: the wheel}
The bindings are packaged as a wheel, one per Python version (Chapter~\ref{ch:packaging} explains wheels):
\begin{lstlisting}[language=bash]
cpp/scripts/build_wheels.sh 3.14        # the Python version of nabia_python/.venv (no argument: 3.10-3.14)
NABIA_WHEEL_BUILD_TYPE=RelWithDebInfo cpp/scripts/build_wheels.sh 3.14   # with debug symbols
cpp\scripts\build_wheels.ps1 -Versions 3.12   # Windows (PowerShell): cibuildwheel, win_amd64
\end{lstlisting}
Each wheel is built in isolation, installed in a clean environment and tested with
\code{nabia\_python/tests/bridge} before it is copied to \code{dist/}. Instead of building, the wheels published
by the CI can be downloaded: \cmd|nabia_python/scripts/fetch_wheels.sh| (the newest successful run, or
\cmd|fetch_wheels.sh release v0.1.0|).

\section{Step 3: the Python project}
\begin{lstlisting}[language=bash]
nabia_python/scripts/setup_env.sh       # creates nabia_python/.venv: the wheel from dist/ + nabia (editable)
cd nabia_python
.venv/bin/python -m pytest              # the bridge tests, the package tests, the pipeline tests
\end{lstlisting}
\code{setup\_env.sh} (\code{setup\_env.ps1} on Windows) reinstalls the wheel every time, so run it again after
rebuilding the C++ code; Python changes are seen at once (editable install). \code{PYTHON=python3.12} chooses the
interpreter of a new \code{.venv}. In PyCharm, open \code{nabia\_python/} and select \code{.venv/bin/python} as the
interpreter. The whole library is then available from Python (Chapter~\ref{c:sec:bridge}):
\begin{lstlisting}[language=Python]
import nabia_python_bridge as nb
from nabia_python_bridge import frames

rb = nb.risk.RiskBudgeting(nb.risk.RiskBudgetingParams(budgets=[0.5, 0.3, 0.2]), nb.risk.Volatility(cov))
weights = rb.optimize()
reg = frames.regress(fund_returns, factor_returns, min_factors=1, max_factors=3, train_fraction=0.7)
\end{lstlisting}

\section{Step 4: the sentiment risk-budgeting pipeline}
From \code{nabia\_python/} (Part~\ref{part:sentiment} describes the method):
\begin{lstlisting}[language=bash]
.venv/bin/python -m nabia.sentiment_rb                    # long-only, offline (synthetic data, mocked LLM), ~20 s
.venv/bin/python -m nabia.sentiment_rb --show             # also open the figures
.venv/bin/python -m nabia.sentiment_rb --config <(echo '{"budgets": {"mode": "long_short"}}')
.venv/bin/python -m nabia.sentiment_rb --ic-sweep 0 0.1 0.2 0.3 0.5    # alpha against news quality
.venv/bin/python -m nabia.sentiment_rb --out output/my_run             # another output folder
.venv/bin/nabia-sentiment --help                          # the same command, installed as a script
\end{lstlisting}
\begin{itemize}
  \item \textbf{Configuration.} Every setting lives in \code{CONFIG} (\coderef{rb_config}): windows, horizon,
        elastic-net grid, refit schedule, normalization, budget bounds, long-only/long-short, LLM models,
        taxonomy, indicators, assets. \code{--config file.json} deep-merges overrides; long-short mode is
        described in \S\ref{sec:rb-longshort}:
\begin{lstlisting}
{"normalize": {"method": "phi"}, "model": {"refit": "ME"}}
{"budgets": {"mode": "long_short", "short_below": 0.35, "long_above": 0.5, "modulate_long_short": false}}
\end{lstlisting}
  \item \textbf{Live LLM}: \cmd|.venv/bin/pip install -e ".[llm]"|, set \code{ANTHROPIC\_API\_KEY}, add
        \code{--live}: the news of the last \code{llm.live\_days} days is scored by Claude (models in
        \code{llm.triage\_model}, \code{scoring\_model}, \code{digest\_model}); the rest stays mocked.
  \item \textbf{Outputs}, in \code{output/sentiment\_rb} by default:
        \begin{itemize}
          \item tables: \code{comparison}, \code{ic}, \code{llm\_ic}, \code{snapshot}, \code{budgets}, \code{signal},
                \code{selection\_frequency}, \code{ic\_in\_out\_of\_sample}, \code{equity}, \code{news} (CSV);
          \item the LLM audit log \code{llm\_audit.jsonl};
          \item figures \code{portfolio.png}, \code{budgets.png}, \code{model.png}; for a sweep,
                \code{sweep.png} and \code{sweep.csv}.
        \end{itemize}
\end{itemize}

\section{Real data}\label{sec:guide-real}
\begin{lstlisting}[language=bash]
.venv/bin/python -m nabia.sentiment_rb --data real --out output/real          # real prices, example news, offline
.venv/bin/python -m nabia.sentiment_rb --data real --news my_news.json         # your own news file
.venv/bin/python -m nabia.sentiment_rb.data.example_news --ic 0.3             # regenerate the example news
.venv/bin/python -m nabia.sentiment_rb --data real --live --batch --config <(echo '{"llm": {"live_days": 0}}')
.venv/bin/python -m nabia.news fetch --no-rss --gkg --since 2024-01-01    # news history -> data/news.json
.venv/bin/python -m nabia.news fetch                                       # today's RSS items (run daily)
.venv/bin/python -m nabia.news score                                     # local Ollama -> data/news_scores.json
.venv/bin/python -m nabia.sentiment_rb --data real --scores data/news_scores.json
\end{lstlisting}
Prices and indicators come from Yahoo Finance and FRED and are cached in \code{data/cache}; news comes from a JSON
file of articles; live scoring caches every response in \code{data/llm\_cache.jsonl}. Chapter~\ref{ch:rb-real}
describes the three parts.

\section{Run configurations and the weights script}\label{sec:guide-weights}
The settings are in three parts (\coderef{rb_config}). \code{CONFIG} holds the methodology every run shares:
the factor tree, assets, features, windows, model, budgets, portfolio and LLM. \code{SYNTHETIC} and \code{REAL}
hold the data settings of each kind of run. A run's JSON file chooses one with \code{data.source} and overrides any
key: \code{configs/synthetic.json} (generated market and news, the oracle skill) and \code{configs/real.json}
(Yahoo and FRED prices, the news and scores files, the LLM provider).
\begin{lstlisting}[language=bash]
.venv/bin/python -m nabia.sentiment_rb --config configs/synthetic.json --out output/syn   # full report and figures
.venv/bin/python -m nabia.sentiment_rb --config configs/real.json --out output/real
.venv/bin/python scripts/sentiment_weights.py configs/real.json                           # weights only
\end{lstlisting}
\code{scripts/sentiment\_weights.py} takes a single argument, the configuration, and computes the signal, the budgets,
the directions and the weights over a lookback window (\code{backtest.lookback\_days}, ending at
\code{backtest.as\_of}, by default the last date). With \code{data.scores\_file} it reads the LLM scores. Without one,
it scores \code{data.news\_file} with \code{llm.provider}: \code{mock} (offline) or \code{ollama} or
\code{anthropic} (live). An \code{as\_of} date also cuts the data there, so the position uses nothing that came later.
It writes the following to \code{backtest.out}:
\begin{itemize}
  \item \code{signal.csv}: the daily $s$;
  \item \code{rebalances.csv}: per rebalance and asset, $s$, the direction, the budget, the weight and the risk
        contribution;
  \item \code{current.csv}: the position to hold from the as-of date;
  \item \code{summary.json}: that position and the window's performance against plain ERC.
\end{itemize}

\section{The staged real-data pipeline}\label{sec:guide-stages}
Five scripts run the real-data pipeline one step at a time, each from the configuration file alone. Each step
writes a file that the later steps read, so it runs once (\coderef{rb_stages}):
\begin{lstlisting}[language=bash]
.venv/bin/python scripts/pipeline/collect_news_1.py configs/real.json      # articles -> data.news_file
.venv/bin/python scripts/pipeline/digest_news_2.py configs/real.json       # LLM daily digests -> data.digests_file
.venv/bin/python scripts/pipeline/score_news_3.py configs/real.json        # LLM article scores -> data.scores_file
.venv/bin/python scripts/pipeline/calibrate_betas_4.py configs/real.json   # elastic-net betas -> data.betas_file
.venv/bin/python scripts/pipeline/backtest_5.py configs/real.json          # portfolio with the stored betas
\end{lstlisting}
\begin{enumerate}
  \item \textbf{Collect} (section \code{collect}): RSS items and GDELT GKG history, deduplicated, appended to the
        news file. Run it again to add new articles.
  \item \textbf{Digest}: the LLM's macro digest of every day (Prompt C), with \code{llm.provider}.
  \item \textbf{Score}: triage and node scores of every article (Prompts A and B). Steps 2 and 3 save their file every
        25 items and cache every response in \code{llm.cache}, so an interrupted run resumes.
  \item \textbf{Calibrate}: the purged walk-forward elastic net of every asset (Chapter~\ref{ch:rb-model}). The betas
        of every refit date are stored with the asset's candidate features and the model settings
        (\coderef{rb_fit_betas}).
  \item \textbf{Backtest}: each refit's stored betas give $S$ and $s$ from the next day to the next refit
        (\coderef{rb_apply_betas}). The budgets, directions, backtest against ERC, report and positions follow. Betas
        whose features are no longer the run's candidates are refused.
\end{enumerate}
Splitting the calibration from the scoring gives exactly the results of the one-shot run, as a test checks.

\section{Debugging the C++ code from Python}\label{sec:guide-debug}
\begin{enumerate}
  \item Build a wheel with debug symbols (\code{NABIA\_WHEEL\_BUILD\_TYPE=RelWithDebInfo}) and run
        \code{setup\_env.sh}; or, without a wheel, \cmd|cmake --preset python-bridge && cmake --build --preset python-bridge|
        in \code{cpp/}, which writes the modules next to the package sources, and run the script with
        \code{PYTHONPATH=<repository>/cpp/bindings}.
  \item Start the script in PyCharm, stopped at a breakpoint before the C++ call.
  \item In CLion: Run $\to$ Attach to Process, the Python process. Breakpoints in \code{cpp/src} are hit.
\end{enumerate}

\section{Wheels for every platform and releases}
\begin{lstlisting}[language=bash]
gh workflow run wheels.yml                          # Linux x86_64/aarch64, macOS arm64/x86_64, Windows AMD64
gh run list --workflow wheels.yml                   # follow it
gh run download <run-id> -n all-wheels -D dist      # every wheel and the sdist
git tag v0.2.0 && git push origin v0.2.0            # a GitHub Release with all of them attached
\end{lstlisting}
Pushes to \code{main} and pull requests that change \code{cpp/} also start the workflow. Chapter~\ref{ch:packaging}
explains every job and how to write new ones.

\section{Building this document}
\begin{lstlisting}[language=bash]
doc/build.sh                 # doc/nabia.pdf: extracts the excerpts from the current code, runs pdflatex 3 times
doc/build.sh --figures       # first reruns the pipeline (long-only, long-short, both sweeps) for the figures/tables
\end{lstlisting}
The excerpts, their file:line locations and the packaging file pieces are extracted from the repository by
\code{doc/extract\_snippets.py}; a function that no longer exists stops the build. The figures and result tables
of Part~\ref{part:sentiment} are copied from \code{nabia\_python/output} into \code{doc/figures} by
\code{--figures}.
